\section{Background and Related Work}
\label{sec:background}

\subsection{Search-Based Software Testing and EvoSuite}
\label{subsec:evosuite}
Automated unit test generation has been extensively studied for over two decades~\cite{yang2025requirements}. The dominant traditional paradigm for Java is Search-Based Software Testing (SBST), exemplified by \textbf{EvoSuite}~\cite{fraser2011evosuite}. EvoSuite formulates test suite generation as a many-objective optimization problem, applying evolutionary algorithms to maximize structural coverage criteria such as line, branch, and mutation coverage. Despite its effectiveness at reaching deep control-flow branches, SBST exhibits two well-documented limitations in industrial practice. First, synthetically evolved test inputs and variable names often lack semantic naturalness, making the resulting JUnit tests difficult for developers to comprehend and maintain~\cite{diop2025scenario, pan2025aster}. Second, because EvoSuite has no independent understanding of intended program semantics, it synthesizes \emph{regression assertions} by executing the current SUT and recording observed return values as expected outputs~\cite{fraser2011evosuite, tang2024chatgpt}. Consequently, SBST tests cannot distinguish intended behavior from latent implementation quirks, motivating the exploration of Large Language Models capable of reasoning over code semantics.

\subsection{LLM-Based Unit Test Generation}
\label{subsec:llm_test_gen}
Pre-trained and instruction-tuned LLMs have recently been adopted to generate unit tests directly from focal source code and documentation~\cite{schafer2024empirical, siddiq2024using, yang2024evaluation}. Siddiq et al.~\cite{siddiq2024using} evaluated Codex, GPT-3.5-Turbo, and StarCoder on HumanEval-Java and the EvoSuite SF110 benchmark, finding that while LLMs achieve over 80\% coverage on simple HumanEval functions, their compilation rates drop precipitously to 2.7\%--12.7\% on complex SF110 classes due to hallucinated APIs and missing context. Similarly, Yang et al.~\cite{yang2024evaluation} evaluated open-source and proprietary LLMs across 17 Defects4J projects, reporting that prompt formulation and focal context strongly influence Compilation Success Rate (CSR), and that compilation failures prevent defect detection in 87.13\% of studied cases.

At an industrial scale, Alshahwan et al.~\cite{alshahwan2024automated} introduced TestGen-LLM at Meta, integrating LLM test generation with strict automated filters that discard non-compiling or non-passing tests before human review. Pan et al.~\cite{pan2025aster} proposed ASTER, combining static program analysis with compilation and runtime remediation across Java SE/EE and Python applications. More recently, Kumar et al.~\cite{kumar2025aipowered} conducted a comparative case study involving GPT-4o, Gemini, and Claude-3.5 on HumanEval-Java, demonstrating that frontier models achieve high compilation success ($>99\%$) on function-level tasks, yet still experience runtime and assertion failures ranging from 9.78\% (Gemini) to 15.28\% (GPT-4o). However, these multi-model comparisons remain largely confined to standalone functions rather than stateful object-oriented classes.

\subsection{The Test Oracle Problem and Prompt Refinement}
\label{subsec:oracle_prompt_refinement}
In software testing, a \emph{test oracle} is the mechanism that determines whether a test execution passes or fails~\cite{yang2025requirements}. In JUnit tests, the oracle is encoded in explicit assertions (e.g., \texttt{assertEquals}, \texttt{assertTrue}, \texttt{assertNull}) or expected exception declarations. Even when an LLM generates syntactically valid Java code, it frequently suffers from \emph{oracle hallucination}: asserting an expected value that contradicts the actual execution semantics of the SUT~\cite{yuan2024evaluating, konstantinou2025yate, ouyang2024empirical}. Recent work on LLM-as-a-Judge further confirms that LLMs frequently misjudge code correctness when relying on superficial reasoning, establishing actual compiler and JUnit test execution as the indispensable ground-truth oracle~\cite{ou2025effectiveness}.

To mitigate generation failures, researchers have explored two primary avenues: \emph{iterative repair loops} and \emph{structured prompt engineering}. In the repair paradigm, Yuan et al.~\cite{yuan2024evaluating} proposed ChatTester, which leverages iterative compiler feedback to raise compilation success from 39.0\% to 73.3\% and execution success from 22.3\% to 41.0\% on Java focal methods. Gu et al.~\cite{gu2024testart} introduced TestART, co-evolving generation and template-based repair across up to four iterations on Defects4J. Konstantinou et al.~\cite{konstantinou2025yate} presented YATE, decomposing test repair into compilation fixing and oracle fixing across six Java projects. However, Singh et al.~\cite{singh2026qualitative} recently demonstrated that multi-round repair loops frequently introduce undesirable \emph{anti-patterns}---such as breaking encapsulation via reflection or mutating assertions merely to force a pass---in 31 out of 117 inspected repair pairs.

Complementing post-execution repair, structured prompting seeks to prevent errors during generation itself. Ryan et al.~\cite{ryan2024codeaware} developed SymPrompt, embedding path constraints into multi-stage prompts to guide path-sensitive test generation. Wang et al.~\cite{wang2024hits} proposed HITS, applying method slicing to decompose complex Java methods with high cyclomatic complexity. Diop et al.~\cite{diop2025scenario} and Chudic and {\c{C}}al{\i}kl{\i}~\cite{chudic2026automated} investigated two-step scenario prompting and retrieval-based few-shot prompting, while Paul et al.~\cite{paul2025impact} showed that structured prompt programming must be carefully controlled to avoid confounding model comparisons. Building on these insights, our \textbf{Oracle-Aware Prompt Refinement (OAPR)} approach targets semantic oracle accuracy directly at the prompt level: rather than feeding back runtime stack traces that might tempt the model to trivialize assertions~\cite{singh2026qualitative}, OAPR enforces a systematic, source-first trace of preconditions, state transitions, and branch-level return semantics.

\subsection{Class-Level Evaluation: ClassEval and ClassEval-T}
\label{subsec:classeval}
A major limitation of early code and test generation benchmarks such as HumanEval~\cite{chen2021evaluating} is their exclusive focus on isolated, stateless functions averaging only 11.5 lines of code (LOC). To evaluate LLMs under realistic object-oriented conditions, Du et al.~\cite{du2024classeval} introduced \textbf{ClassEval} at ICSE 2024---the first manually crafted class-level Python benchmark comprising 100 classes designed under the \emph{Design by Contract} principle ($\sim$500 person-hours). Unlike HumanEval, ClassEval classes encapsulate mutable instance fields, constructors, and coupled methods.

Subsequently, Xue et al.~\cite{xue2025classevalt} presented \textbf{ClassEval-T} at ISSTA 2025, investing 360 person-hours of manual expert translation to port ClassEval into parallel \textbf{Java} and C++ benchmarks. Excluding six tasks that relied on Python-exclusive libraries with no Java or C++ equivalents, ClassEval-T establishes 94 parallel Java classes accompanied by high-coverage reference test suites achieving 99.7\% statement coverage and 98.2\% branch coverage. Because 76.2\% of the methods in ClassEval-T exhibit internal dependencies on instance fields or sibling methods (averaging 66.7 LOC per class), ClassEval-T provides an ideal, contamination-resistant testbed for investigating class-level test oracle failures and evaluating OAPR across ChatGPT and Gemini.
